
\documentclass[border=4pt]{standalone}
\usepackage{polyglossia}
\setdefaultlanguage{russian}
\setotherlanguage{english}
\usepackage{fontspec}
\IfFontExistsTF{CMU Serif}{\setmainfont{CMU Serif}}{\setmainfont{Liberation Serif}}
\IfFontExistsTF{CMU Sans Serif}{\setsansfont{CMU Sans Serif}}{\setsansfont{Liberation Sans}}
\IfFontExistsTF{CMU Typewriter Text}{\setmonofont{CMU Typewriter Text}}{\setmonofont{DejaVu Sans Mono}}
\usepackage{amsmath,amssymb,mathtools,bm}
\usepackage{xcolor}
\definecolor{accent}{HTML}{1F4E79}
\definecolor{accent2}{HTML}{B85450}
\definecolor{shade}{HTML}{F4F1EA}
\colorlet{prosvBlue}{accent}
\colorlet{prosvRed}{accent2}
\colorlet{prosvLight}{accent!12}
\colorlet{prosvCream}{shade}
\definecolor{prosvGold}{HTML}{E8B647}
\definecolor{prosvGreen}{HTML}{6A9F58}
\colorlet{prosvGray}{black!55}
\usepackage{tikz}
\usetikzlibrary{calc, arrows.meta, decorations.pathreplacing,
                positioning, shapes.geometric, shapes.misc, fit,
                backgrounds}
\usepackage{pgfplots}
\pgfplotsset{compat=1.18}
\newcommand{\R}{\mathbb{R}}
\newcommand{\N}{\mathbb{N}}
\newcommand{\Z}{\mathbb{Z}}
\newcommand{\eps}{\varepsilon}
\newcommand{\norm}[1]{\left\|#1\right\|}
\newcommand{\abs}[1]{\left|#1\right|}
\newcommand{\NOD}{\gcd}
\newcommand{\Eul}{\varphi}
\newcommand{\Zn}[1]{\mathbb{Z}_{#1}}
\newcommand{\modn}[1]{\,(\mathrm{mod}\,#1)}
\newcommand{\term}[1]{\textcolor{accent2}{\textbf{#1}}}
\newcommand{\engterm}[1]{\textit{#1}}
\DeclareMathOperator*{\argmin}{arg\,min}
\DeclareMathOperator*{\argmax}{arg\,max}
\begin{document}

\begin{tikzpicture}[
  every node/.style={font=\footnotesize},
  cell/.style={draw=prosvBlue, fill=prosvLight, minimum width=0.7cm, minimum height=0.7cm},
  hl/.style={draw=prosvRed, fill=prosvCream, minimum width=0.7cm, minimum height=0.7cm, thick}
]
  \node[font=\small, prosvBlue] at (-2, 0.7) {Счётчики:};
  \node[cell] at (-1.3, 0) {3};
  \node[cell] at (-0.6, 0) {1};
  \node[hl]   at (0.1, 0)  {7};
  \node[cell] at (0.8, 0)  {2};
  \node[cell] at (1.5, 0)  {4};
  \node[cell] at (2.2, 0)  {0};
  \node[cell] at (2.9, 0)  {3};
  \node[cell] at (3.6, 0)  {1};
  \node[cell] at (4.3, 0)  {2};
  \node[cell] at (5.0, 0)  {5};
  \node at (5.7, 0)   {$\dots$};

  \node[align=left, font=\footnotesize, prosvRed] at (8, 0) {\textit{редкий «хвост»} ---\\ скорее всего \\ много элементов};
  \draw[->, prosvRed] (8, -0.5) -- (0.4, -0.3);

  \node[font=\footnotesize\itshape, prosvBlue, align=center] at (3, -1.2)
    {усреднение по всем счётчикам с помощью гармонического среднего};
\end{tikzpicture}
\end{document}
